% ============================================================
%  CUADERNILLO DE ESTUDIO — CLASE 5
%  Seguridad de la Información y de Sistemas
%  Lic. en Gestión de Negocios Digitales — UCA
%  Compilar con: pdflatex (dos pasadas para el índice)
% ============================================================

\documentclass[12pt,a4paper]{article}

% ---------- Paquetes ----------
\usepackage[utf8]{inputenc}
\usepackage[spanish]{babel}
\usepackage[T1]{fontenc}
\usepackage{lmodern}
\usepackage[margin=2.2cm,headheight=14pt]{geometry}
\usepackage{amsmath}
\usepackage{amssymb}
\usepackage{enumitem}
\usepackage{booktabs}
\usepackage{array}
\usepackage{tabularx}
\usepackage[table]{xcolor} % Necesario para \cellcolor en tablas
\usepackage{tcolorbox}
\usepackage{graphicx}
\usepackage{fancyhdr}
\usepackage{titlesec}
\usepackage{hyperref}
\usepackage{multicol}
\usepackage{pifont}
\usepackage{wasysym}
\usepackage{tikz}
\usetikzlibrary{positioning,arrows.meta,shapes.geometric}

% ---------- Colores ----------
\definecolor{azulUCA}{HTML}{003366}
\definecolor{doradoUCA}{HTML}{B8860B}
\definecolor{grisClaro}{HTML}{F2F2F2}
\definecolor{rojoAlerta}{HTML}{C0392B}
\definecolor{verdeOK}{HTML}{27AE60}
\definecolor{azulCielo}{HTML}{2980B9}
\definecolor{naranjaAmenaza}{HTML}{D35400}

% Colores para la Matriz de Riesgos (Mapa de Calor)
\definecolor{riskLow}{HTML}{A9DFBF}   % Verde
\definecolor{riskMed}{HTML}{F9E79F}   % Amarillo
\definecolor{riskHigh}{HTML}{F5CBA7}  % Naranja
\definecolor{riskCrit}{HTML}{F1948A}  % Rojo

% ---------- tcolorbox estilos ----------
\tcbuselibrary{skins,breakable}

\newtcolorbox{cajaTitulo}[1][]{
  colback=azulUCA, colframe=azulUCA, coltext=white,
  fonttitle=\Large\bfseries, arc=4mm, #1
}

\newtcolorbox{cajaDefinicion}{
  colback=azulCielo!10, colframe=azulCielo,
  fonttitle=\bfseries, title=Definici\'on, arc=3mm, breakable
}

\newtcolorbox{cajaProceso}{
  colback=doradoUCA!10, colframe=doradoUCA,
  fonttitle=\bfseries, title=El Proceso, arc=3mm, breakable
}

\newtcolorbox{cajaActividad}{
  colback=grisClaro, colframe=azulUCA,
  fonttitle=\bfseries\color{azulUCA},
  title={\ding{45}\; Actividad Pr\'actica}, arc=3mm, breakable
}

\newtcolorbox{cajaNota}{
  colback=rojoAlerta!8, colframe=rojoAlerta,
  fonttitle=\bfseries\color{rojoAlerta}, title=Nota de Negocio, arc=3mm, breakable
}

% ---------- Header / Footer ----------
\pagestyle{fancy}
\fancyhf{}
\fancyhead[L]{\small\textcolor{azulUCA}{Seguridad de la Informaci\'on y de Sistemas}}
\fancyhead[R]{\small\textcolor{azulUCA}{Cuadernillo --- Clase 5}}
\fancyfoot[C]{\thepage}
\fancyfoot[L]{\small\textcolor{gray}{UCA}}
\renewcommand{\headrulewidth}{0.4pt}
\renewcommand{\footrulewidth}{0.4pt}

% ---------- Títulos de sección ----------
\titleformat{\section}
  {\color{azulUCA}\Large\bfseries}
  {\thesection.}{0.6em}{}
\titleformat{\subsection}
  {\color{azulCielo}\large\bfseries}
  {\thesubsection}{0.6em}{}

% ---------- Enlaces ----------
\hypersetup{colorlinks=true, linkcolor=azulUCA, urlcolor=azulCielo}

% ---------- Checkbox ----------
\newcommand{\checkboxVacio}{$\square$\;}
\newcommand{\checkMark}{\ding{51}}

% ============================================================
\begin{document}

% ============================================================
%  PORTADA
% ============================================================
\begin{titlepage}
\centering
\vspace*{1cm}

\begin{cajaTitulo}[width=\textwidth, halign=center]
  SEGURIDAD DE LA INFORMACI\'ON Y DE SISTEMAS\\[4pt]
  {\large Licenciatura en Gesti\'on de Negocios Digitales --- UCA}
\end{cajaTitulo}

\vspace{1.5cm}

{\Huge\bfseries\color{azulUCA} CUADERNILLO DE ESTUDIO}\\[8pt]
{\LARGE\color{doradoUCA} CLASE 5}\\[12pt]
{\Large\itshape ``Gesti\'on de riesgos I:\\[2pt] Metodología y Matriz P$\times$I''}\\[1.5cm]

\begin{tabular}{rl}
\textbf{Profesor:}   & Mg.\ Ing.\ Rodrigo Atilio Elgueta \\[4pt]
\textbf{Unidad:}     & Unidad 2 \\[4pt]
\end{tabular}

\vfill

\begin{tcolorbox}[colback=grisClaro, colframe=azulUCA, arc=3mm, width=0.9\textwidth]
  \centering
  {\bfseries Datos del/la estudiante}\\[10pt]
  Nombre y apellido: \underline{\hspace{5cm}}\\[10pt]
  Grupo N.\textdegree: \underline{\hspace{5cm}}\\[10pt]
  Negocio Digital (TP): \underline{\hspace{5cm}}\\[10pt]
  Fecha: \underline{\hspace{5cm}}
\end{tcolorbox}

\vspace{1cm}
{\small\textcolor{gray}{Material de uso exclusivo para la cursada.}}
\end{titlepage}

% ============================================================
%  ÍNDICE
% ============================================================
\tableofcontents
\newpage

% ============================================================
%  SECCIÓN 1 — INTRODUCCIÓN
% ============================================================
\section{¿Qu\'e es gestionar riesgos en un negocio digital?}

En las clases anteriores identificamos \textbf{activos}, \textbf{amenazas} y \textbf{vulnerabilidades}. Hoy damos el siguiente paso: ¿qu\'e hacemos con todo eso? 

\begin{cajaDefinicion}
La \textbf{Gesti\'on de Riesgos} es el proceso de identificar, analizar y evaluar los riesgos para luego decidir c\'omo tratarlos, optimizando los recursos del negocio.
\end{cajaDefinicion}

\begin{cajaNota}
\textbf{El objetivo NO es el ``Riesgo Cero''.} Eliminar todos los riesgos es t\'ecnicamente imposible y econ\'omicamente inviable (arruinaría la usabilidad y el presupuesto). El objetivo es \textbf{administrar el riesgo} hasta un nivel aceptable para el negocio, invirtiendo en controles solo donde el costo-beneficio lo justifique.
\end{cajaNota}

\section{El Ciclo de Gesti\'on de Riesgos}

El proceso es iterativo y c\'iclico. Consta de 4 etapas principales:

\begin{cajaProceso}
\begin{enumerate}[leftmargin=2em, itemsep=8pt]
    \item \textbf{Identificaci\'on:} ¿Qu\'e puede salir mal? (Listar activos, amenazas y vulnerabilidades).
    \item \textbf{An\'alisis:} ¿Qu\'e tan grave ser\'ia y qu\'e tan probable es que ocurra? (Estimar Probabilidad e Impacto).
    \item \textbf{Evaluaci\'on:} ¿Qu\'e riesgos son prioritarios? (Usar la Matriz P$\times$I para rankearlos).
    \item \textbf{Tratamiento:} ¿Qu\'e vamos a hacer al respecto? (Mitigar, Transferir, Aceptar o Evitar).
\end{enumerate}
\end{cajaProceso}

\newpage
% ============================================================
%  SECCIÓN 2 — LA MATRIZ P x I
% ============================================================
\section{La Matriz de Probabilidad e Impacto (P$\times$I)}

Para no tratar todos los riesgos por igual, los cruzamos en una matriz cualitativa. Esto nos da un \textbf{Mapa de Calor} que dicta la urgencia.

\subsection{Escalas de medici\'on}
\begin{itemize}[leftmargin=2em]
    \item \textbf{Probabilidad (P):} ¿Qu\'e tan frecuente o f\'acil es que la amenaza explote la vulnerabilidad? (Baja, Media, Alta).
    \item \textbf{Impacto (I):} Si ocurre, ¿cu\'anto da\~no causa al negocio (financiero, legal, reputacional, operativo)? (Bajo, Medio, Alto).
\end{itemize}

\subsection{Mapa de Calor (Heatmap)}

\begin{center}
\renewcommand{\arraystretch}{1.5}
\begin{tabular}{|c|>{\columncolor{grisClaro}}c|>{\columncolor{grisClaro}}c|>{\columncolor{grisClaro}}c|}
\hline
\rowcolor{azulUCA!20}
\textbf{Prob. $\downarrow$ \quad Impacto $\rightarrow$} & \textbf{Bajo} & \textbf{Medio} & \textbf{Alto} \\
\hline
\textbf{Alta} & \cellcolor{riskMed} \textbf{MEDIO} & \cellcolor{riskHigh} \textbf{ALTO} & \cellcolor{riskCrit} \textbf{CR\'ITICO} \\
\hline
\textbf{Media} & \cellcolor{riskLow} \textbf{BAJO} & \cellcolor{riskMed} \textbf{MEDIO} & \cellcolor{riskHigh} \textbf{ALTO} \\
\hline
\textbf{Baja} & \cellcolor{riskLow} \textbf{BAJO} & \cellcolor{riskLow} \textbf{BAJO} & \cellcolor{riskMed} \textbf{MEDIO} \\
\hline
\end{tabular}
\end{center}

\vspace{8pt}
\textbf{¿C\'omo se lee?}
\begin{itemize}[leftmargin=2em]
    \item \textcolor{red!80!black}{\textbf{Cr\'itico / Alto:}} Requieren tratamiento inmediato. No se pueden dejar pasar.
    \item \textcolor{orange!80!black}{\textbf{Medio:}} Requieren un plan de acci\'on a mediano plazo y monitoreo.
    \item \textcolor{green!60!black}{\textbf{Bajo:}} Se aceptan o se monitorean peri\'odicamente.
\end{itemize}

\newpage
% ============================================================
%  SECCIÓN 3 — TRATAMIENTO DE RIESGOS
% ============================================================
\section{Opciones de Tratamiento (Las 4 T)}

Una vez evaluado el riesgo, la direcci\'on del negocio debe decidir qu\'e hacer. Existen 4 caminos cl\'asicos:

\renewcommand{\arraystretch}{1.6}
\begin{tabularx}{\textwidth}{|>{\bfseries\color{azulUCA}}p{3cm}|X|p{4cm}|}
\hline
\rowcolor{grisClaro}
\textbf{Estrategia} & \textbf{Descripci\'on} & \textbf{Ejemplo en Negocio Digital} \\
\hline
\textbf{Mitigar} (Reducir) & Aplicar controles para bajar la Probabilidad o el Impacto. & Implementar MFA (doble factor) para reducir la probabilidad de robo de credenciales. \\
\hline
\textbf{Transferir} (Compartir) & Pasar el impacto financiero a un tercero. & Contratar un \textbf{seguro de ciber-riesgo} o tercerizar el hosting a un proveedor con SLA. \\
\hline
\textbf{Aceptar} (Tolerar) & Asumir el riesgo porque el control es m\'as caro que el da\~no potencial. & Aceptar que el servidor se caiga 10 min al a\~no porque un backup redundante cuesta u\$s 50k y la p\'erdida es de u\$s 500. \\
\hline
\textbf{Evitar} (Eliminar) & Dejar de realizar la actividad que genera el riesgo. & Decidir NO almacenar datos de tarjetas de cr\'edito propios y usar siempre pasarelas como Mercado Pago o Stripe. \\
\hline
\end{tabularx}

\begin{cajaNota}
\textbf{El BIA (Business Impact Analysis):} Antes de decidir, hay que entender el impacto real en el negocio. Si el sistema de facturaci\'on cae 2 horas, ¿cu\'anto dinero se pierde? Eso dicta cu\'anto estamos dispuestos a invertir en prevenirlo.
\end{cajaNota}

\newpage
% ============================================================
%  SECCIÓN 4 — ACTIVIDAD PRÁCTICA
% ============================================================
\section{Taller Grupal: Matriz para el TP Integrador}

\begin{cajaActividad}[title={Construcci\'on de la Matriz de Riesgos}]
En sus grupos del Trabajo Pr\'actico Integrador, identifiquen \textbf{3 riesgos clave} para el negocio digital que eligieron. 

\textbf{Pasos:}
\begin{enumerate}
    \item Describan el riesgo (Amenaza + Vulnerabilidad).
    \item Identifiquen el Activo afectado.
    \item As\'ignenle un valor de Probabilidad (B/M/A) e Impacto (B/M/A).
    \item Determinen el Nivel de Riesgo seg\'un el Mapa de Calor.
    \item Propongan un Control y la Estrategia (Mitigar, Transferir, etc.).
\end{enumerate}
\end{cajaActividad}

\vspace{8pt}
\renewcommand{\arraystretch}{1.8}
\begin{tabularx}{\textwidth}{|X|}
\hline
\textbf{Riesgo 1:} \hrulefill \\
\textbf{Activo:} \hrulefill \quad \textbf{Prob:} [ B / M / A ] \quad \textbf{Imp:} [ B / M / A ] \quad \textbf{Nivel:} \hrulefill \\
\textbf{Control propuesto:} \hrulefill \\
\hline
\end{tabularx}

\vspace{8pt}
\begin{tabularx}{\textwidth}{|X|}
\hline
\textbf{Riesgo 2:} \hrulefill \\
\textbf{Activo:} \hrulefill \quad \textbf{Prob:} [ B / M / A ] \quad \textbf{Imp:} [ B / M / A ] \quad \textbf{Nivel:} \hrulefill \\
\textbf{Control propuesto:} \hrulefill \\
\hline
\end{tabularx}

\vspace{8pt}
\begin{tabularx}{\textwidth}{|X|}
\hline
\textbf{Riesgo 3:} \hrulefill \\
\textbf{Activo:} \hrulefill \quad \textbf{Prob:} [ B / M / A ] \quad \textbf{Imp:} [ B / M / A ] \quad \textbf{Nivel:} \hrulefill \\
\textbf{Control propuesto:} \hrulefill \\
\hline
\end{tabularx}

\vspace{12pt}
\begin{cajaNota}
\textbf{Nota para el TP:} Esta tabla es el borrador directo del Entregable 2 de su Trabajo Pr\'actico Integrador. Gu\'ardenla y s\'aquenle foto/escan\'eenla para su informe.
\end{cajaNota}

\newpage
% ============================================================
%  SECCIÓN 5 — CUESTIONARIO 2
% ============================================================
\section{Cuestionario 2 de autoevaluaci\'on}

\begin{cajaNota}[title={Instrucciones}]
Este cuestionario se realiza al cierre de la clase. \textbf{No es eliminatorio.} Usalo para medir cu\'anto entendiste antes del Parcial 1 (Clase 6). Las respuestas correctas est\'an al final.
\end{cajaNota}

\vspace{6pt}

\textbf{1.} El proceso de gesti\'on de riesgos sigue, en orden:
\begin{enumerate}[label=\alph*), leftmargin=2.5em]
  \item Tratamiento, identificaci\'on, an\'alisis, evaluaci\'on
  \item Identificaci\'on, an\'alisis, evaluaci\'on, tratamiento
  \item Evaluaci\'on, identificaci\'on, tratamiento, an\'alisis
  \item No sigue un orden espec\'ifico
\end{enumerate}

\textbf{2.} En una matriz de probabilidad e impacto, un riesgo con alta probabilidad y alto impacto debe:
\begin{enumerate}[label=\alph*), leftmargin=2.5em]
  \item Ignorarse por ser poco com\'un
  \item Priorizarse para tratamiento inmediato
  \item Transferirse siempre a un seguro
  \item Aceptarse sin m\'as
\end{enumerate}

\textbf{3.} BIA significa:
\begin{enumerate}[label=\alph*), leftmargin=2.5em]
  \item Business Impact Analysis (An\'alisis de Impacto al Negocio)
  \item Basic Information Access
  \item Business Information Audit
  \item Backup and Incident Assessment
\end{enumerate}

\textbf{4.} Un ejemplo de control t\'ecnico es:
\begin{enumerate}[label=\alph*), leftmargin=2.5em]
  \item Una pol\'itica firmada por empleados
  \item Un firewall o un sistema de autenticaci\'on multifactor (MFA)
  \item Una capacitaci\'on anual
  \item Un cerco perimetral
\end{enumerate}

\textbf{5.} Un ejemplo de control humano/organizativo es:
\begin{enumerate}[label=\alph*), leftmargin=2.5em]
  \item Antivirus
  \item Capacitaci\'on en concientizaci\'on de seguridad
  \item Cifrado de disco
  \item C\'amaras de seguridad
\end{enumerate}

\textbf{6.} ISO/IEC 27001 es principalmente:
\begin{enumerate}[label=\alph*), leftmargin=2.5em]
  \item Una ley argentina de protecci\'on de datos
  \item Una norma internacional sobre sistemas de gesti\'on de seguridad de la informaci\'on
  \item Un lenguaje de programaci\'on seguro
  \item Un tipo de malware
\end{enumerate}

\textbf{7.} ``Tratar'' un riesgo puede significar (eleg\'i todas las correctas):
\begin{enumerate}[label=\alph*), leftmargin=2.5em]
  \item Evitarlo
  \item Mitigarlo
  \item Transferirlo
  \item Aceptarlo
\end{enumerate}

\textbf{8.} Verdadero o Falso: la gesti\'on de riesgos es un proceso que se hace una sola vez y no se repite.
\vspace{4pt}
\checkboxVacio Verdadero \qquad \checkboxVacio Falso

% --- Clave de respuestas ---
\vspace{12pt}
\begin{tcolorbox}[colback=verdeOK!10, colframe=verdeOK,
                  title={\checkMark \; Clave de respuestas}, arc=3mm]
\begin{enumerate}[leftmargin=2em]
  \item \textbf{b)} Identificaci\'on, an\'alisis, evaluaci\'on, tratamiento.
  \item \textbf{b)} Priorizarse para tratamiento inmediato.
  \item \textbf{a)} Business Impact Analysis.
  \item \textbf{b)} Un firewall o un sistema MFA.
  \item \textbf{b)} Capacitaci\'on en concientizaci\'on.
  \item \textbf{b)} Una norma internacional (SGSI).
  \item \textbf{a, b, c, d} (Todas son estrategias v\'alidas de tratamiento).
  \item \textbf{Falso.} Es un proceso continuo y c\'iclico que debe revisarse peri\'odicamente.
\end{enumerate}
\end{tcolorbox}

\newpage
% ============================================================
%  SECCIÓN 6 — GLOSARIO Y NOTAS
% ============================================================
\section{Glosario de la Clase 5}

\renewcommand{\arraystretch}{1.4}
\begin{tabularx}{\textwidth}{>{\bfseries}p{4.5cm}X}
\toprule
\textbf{T\'ermino} & \textbf{Definici\'on en una l\'inea} \\
\midrule
Gesti\'on de Riesgos & Proceso de identificar, analizar, evaluar y tratar riesgos. \\
Matriz P$\times$I & Herramienta visual para cruzar Probabilidad e Impacto. \\
Mapa de Calor & Representaci\'on gr\'afica de la matriz con colores de urgencia. \\
BIA & An\'alisis de Impacto al Negocio (cu\'anto cuesta que algo se detenga). \\
Mitigar & Aplicar controles para reducir la probabilidad o el impacto. \\
Transferir & Traspasar el impacto financiero (ej. seguros, tercerizaci\'on). \\
Aceptar & Asumir el riesgo porque el control es m\'as caro que el da\~no. \\
Evitar & Eliminar la actividad que genera el riesgo. \\
ISO/IEC 27001 & Norma internacional para Sistemas de Gesti\'on de Seguridad (SGSI). \\
\bottomrule
\end{tabularx}

\vspace{1cm}

\section{Espacio para notas y hallazgos}

\begin{tcolorbox}[colback=grisClaro, colframe=gray!50, arc=2mm, height=7cm]
\textcolor{gray}{\itshape Anot\'a aqu\'i dudas sobre la matriz, ideas de controles para tu TP Integrador o conceptos para el Parcial 1.}
\end{tcolorbox}

\vfill
\begin{center}
\small\textcolor{gray}{%
Cuadernillo elaborado para la c\'atedra de Seguridad de la Informaci\'on y de
Sistemas --- Lic.\ en Gesti\'on de Negocios Digitales, UCA.\\
Prohibida su distribuci\'on fuera del \'ambito de la cursada sin autorizaci\'on del docente.}
\end{center}

\end{document}